\chapter{Diagnostic Causal Reasoning}
\label{sec:p19-rationality-7-3}

Diagnosis can ask whether a candidate cause is present, whether it produced
an observed effect, or which fault assumptions make observations consistent
with a system description. Meder and Mayrhofer's chapter~7.3 addresses the
first two questions through probabilistic causal models and experiments
on human judgments. It does not present Reiter's consistency-based diagnosis
as a psychological explanation. The chapter's central argument is that
observed contingency information and an inferred generating causal model
are different representational levels. Their answers may coincide under
one set of assumptions and separate under another.
\cite[chapter~7.3, pp.~449--455]{MederMayrhofer2021DiagnosticCausal}

The treatment below retains the source's progression: elementary statistical
diagnosis; noisy-OR causal diagnosis; parameter and structure uncertainty;
causal attribution; networks with several causes or effects; and the
implications for evaluating rationality. Its mathematical illustrations
are stipulated models, not medical advice. The central empirical record is
the original study used for the chapter's Figure~7.3.1, with its student
sample, fictional patient stimuli and instructed response task.

\section{Select the diagnostic task before selecting a norm}
\label{sec:p19-rat73-task}

\begin{hbdefinition}{def:p19-rat73-queries}{Presence, production and logical compatibility}
For the elementary probabilistic task, \(C,E\) are binary variables.
The lowercase symbols \(c,e\) denote their presence events and
\(\neg c,\neg e\) their absence events. One specified probability law
\(P\) governs each calculation, and every conditioned event has positive
probability. A data table reports joint observations of \(C,E\).

\emph{Predictive probability} is \(P(e\mid c)\): the chance of the effect
given the cause's presence. \emph{Diagnostic probability} is
\(P(c\mid e)\): the chance of the cause's presence given the effect.
Neither quantity alone states that the candidate generated this effect.
\emph{Causal attribution} asks for the probability of a specified production
event. Following the source, write \(c\to e\) for that event, with its
mechanism semantics explicitly supplied below. This arrow is not material
implication and is not itself an intervention operator.

By contrast, the consistency-based task in
\cref{sec:p16-consistency-based-diagnosis} takes a system theory,
observations and normality assumptions for a fixed component set. It
returns inclusion-minimal fault assignments that restore consistency.
Their existence or minimality does not rank their probabilities or identify
an effect-producing mechanism.
\index{diagnostic reasoning!presence versus attribution}
\index{diagnosis!empirical causal reasoning}
\end{hbdefinition}

The source's disease/symptom example is a query template: given the symptom,
estimate whether the disease is present. Its genetic-disposition example
separates that question from whether a disposition produced the symptoms.
These examples motivate the quantities; they are not observational patient
studies. Knowing that a disposition occurs in half the population does not
make it responsible for half the symptoms if it has no causal influence.
\cite[secs.~1--3, pp.~449--452]{MederMayrhofer2021DiagnosticCausal}

\begin{hbexample}{def:p19-rat73-logical-contrast}{Two diagnoses, different probabilities}
Suppose a fictional alarm is generated exactly when at least one of two
faults \(A,C\) is present. A system description contains
\(B\leftrightarrow(A\lor C)\), and the observation is \(B\).
The all-normal assumptions \(\neg A,\neg C\) conflict with that observation.
Either singleton fault assignment restores consistency; both are
inclusion-minimal. Now add the separate probabilistic assumptions that
the faults are independent and \(P(A)=0.1,P(C)=0.4\).
Then \(P(B)=0.46\), \(P(A\mid B)=5/23\), and
\(P(C\mid B)=20/23\), by the calculation in
\cref{thm:p19-rat73-network}. The two probabilities need not sum to one,
because both faults can occur. An inclusion-minimal diagnosis is not an
exclusive probabilistic hypothesis. Conversely, an observation with no
logical contradiction can still change a cause's probability.
\end{hbexample}

\section{From observed contingency to a generating model}
\label{sec:p19-rat73-elementary}

\subsection{The data-level benchmark}

The source's three samples each contain forty fictitious cases:
\[
\begin{array}{c|rrrr|cc}
 &N(c,e)&N(c,\neg e)&N(\neg c,e)&N(\neg c,\neg e)
 &\widehat P(e\mid c)&\widehat P(c\mid e)\\\hline
 1&6&14&2&18&0.3&0.75\\
 2&12&8&4&16&0.6&0.75\\
 3&18&2&6&14&0.9&0.75
\end{array}
\]
Here \(N\) means the number of cases in a cell, not participants, and a
hat means the relative frequency computed from that particular table.
The cause's relative frequency is \(0.5\) throughout. Counting the
effect-positive cases gives \(6/8=12/16=18/24=3/4\).
The change in predictive probability is therefore not a change in the
observed diagnostic proportion. These are the chapter's actual
Figure~7.3.1 inputs, not a newly collected dataset.
\cite[fig.~7.3.1 and eq.~(1), p.~450]{MederMayrhofer2021DiagnosticCausal}

\begin{hbtheorem}{thm:p19-rat73-elementary}{proved=annex:p19.rat73-elementary}{Statistical diagnosis, causal power and attribution}
Let \(0<P(c)<1\), \(P(e)>0\), and put
\(b_c=P(c)\), \(u=P(e\mid c)\), \(v=P(e\mid\neg c)\).
Then
\[
 P(c\mid e)=\frac{u b_c}{u b_c+v(1-b_c)}.
\]
For the additional noisy-OR construction, take mutually independent
binary variables \(C,X,Z\), with presence probabilities
\(b_c,w_c,w_a\), respectively, and define
\(E=(C\land X)\lor Z\). Here \(X\) denotes success of the candidate's
productive mechanism, and \(Z\) amalgamates successful background
production; their parameters lie in \([0,1]\).
Then
\[
\begin{split}
 u&=w_c+w_a-w_cw_a,\qquad v=w_a,\\
 P(e)&=b_cw_c+w_a-b_cw_cw_a,\\
 P(c\mid e)&=
 \frac{b_cw_c+b_cw_a-b_cw_cw_a}
      {b_cw_c+w_a-b_cw_cw_a}.
\end{split}
\]
If \(v<1\), the model identifies
\(w_c=(u-v)/(1-v)\); its generative range requires \(u\geq v\).
Define the production event \(c\to e\) to be \(C\land X\).
With the same positive denominator,
\[
 P(c\to e\mid e)=
 \frac{b_cw_c}{b_cw_c+w_a-b_cw_cw_a}
 \ \leq\ P(c\mid e).
\]
The difference is
\(b_cw_a(1-w_c)/P(e)\). If \(w_c=0\), attribution is zero
while diagnosis equals \(b_c\); if \(w_a=0\), both probabilities
are one whenever \(P(e)>0\).
\end{hbtheorem}

The first equation is a probability identity, applicable without a causal
interpretation. The later equations need a productive mechanism:
independent candidate/background production, no preventive interaction,
stable powers independent of cause occurrence, and no unmodeled production
outside the amalgamated background. Those are the source's power-PC/noisy-OR
assumptions, not consequences of Bayes' rule. An inhibitory contingency
\(u<v\) requires another mechanism; setting a negative number as a
generative probability is not a repair. If \(v=1\), the background always
produces the effect and this inversion cannot identify \(w_c\).
\cite[sec.~2.2, eqs.~(2)--(3), p.~451]{MederMayrhofer2021DiagnosticCausal}
\cite[pp.~279--280]{MederMayrhoferWaldmann2014StructureInduction}

\begin{hbexample}{def:p19-rat73-notes}{The two source notes with exact quantities}
For Sample~1, the plug-in values are \(b_c=1/2,w_a=1/10\) and
\(w_c=(3/10-1/10)/(1-1/10)=2/9\).
Thus \(P(e)=1/5\), the diagnosis is \(3/4\), and the production
attribution is \((1/9)/(1/5)=5/9\), approximately \(0.56\).
The source rounds \(w_c\) to \(0.22\); keeping \(2/9\) makes
the two notes' calculation exact. The causal model and the empirical
Bayes calculation give the same diagnostic value when the parameters
are fitted this way, although they represent the task differently.
The lower attribution discounts cases where the candidate was present
but its mechanism did not produce the effect.
\cite[notes~1--2, pp.~454--455]{MederMayrhofer2021DiagnosticCausal}
\end{hbexample}

Production here can include simultaneous success of target and background.
It does not mean that the target was the sole productive cause or was
necessary for the observed outcome. Counterfactual necessity and a decision
about which intervention to perform need further semantics and loss/cost
assumptions. The finite structural model in
\cref{sec:p16-finite-causal-models} supplies equation-replacement semantics;
a diagnostic conditional alone does not supply those intervention answers.

\subsection{Parameter uncertainty is not structure uncertainty}

A point estimate plugs one fitted parameter tuple into the causal
calculation. Parameter uncertainty instead assigns a distribution to
possible values, updates it with the data, and retains the spread.
Structure uncertainty concerns which generating graph applies.
The source compares \(S_1\), with candidate and background producing
the effect, to \(S_0\), with only background production and no link
from \(C\) to \(E\). Under \(S_0\), the effect provides no diagnostic
evidence about the candidate at fixed parameters: \(P(c\mid e)=b_c\).
This is a two-graph hypothesis space, not a search over every possible
confounding, inhibitory or selection mechanism.
\cite[sec.~2.3, pp.~451--452]{MederMayrhofer2021DiagnosticCausal}

\begin{hbmethod}{def:p19-rat73-structure-method}{The source's structure-induction rating rule}
Let \(D\) be the four observed cell counts, \(\theta_i\) the
parameters for \(S_i\), and \(L_i(D\mid\theta_i)\) the likelihood
of independent cases. \(S_1\) has \((b_c,w_c,w_a)\);
\(S_0\) has \((b_c,w_a)\) and fixes \(w_c=0\).
The four model cell probabilities are
\[
\begin{array}{c|cc}
 &S_0&S_1\\\hline
 c,e&b_cw_a&b_c(w_c+w_a-w_cw_a)\\
 c,\neg e&b_c(1-w_a)&b_c(1-w_c)(1-w_a)\\
 \neg c,e&(1-b_c)w_a&(1-b_c)w_a\\
 \neg c,\neg e&(1-b_c)(1-w_a)&(1-b_c)(1-w_a)
\end{array}
\]
The likelihood is the product of those four probabilities raised to
their respective cell counts, with \(0^0=1\) for zero counts.
A common multinomial coefficient, if the counts rather than an ordered
sample are the observation, cancels in the posterior comparisons.
Choose the source's independent uniform parameter priors
\(\pi_i\) and equal graph priors. Set
\[
\begin{split}
 m_i(D)&=\int L_i(D\mid\theta_i)\pi_i(d\theta_i),\\
 \pi_i(d\theta_i\mid D)&=
       L_i(D\mid\theta_i)\pi_i(d\theta_i)/m_i(D),\\
 \rho_i(D)&=m_i(D)/(m_0(D)+m_1(D)),\\
 g_i(D)&=\int P_{\theta_i,S_i}(c\mid e)
                         \pi_i(d\theta_i\mid D),\\
 g(D)&=\rho_0(D)g_0(D)+\rho_1(D)g_1(D).
\end{split}
\]
Integration is over each graph's own parameter domain; the evidence
normalizers must be positive. The conditional ratio is defined almost
everywhere under these continuous uniform priors, so boundary points with
zero effect probability do not contribute. Other priors must check this
condition anew. The 2014 implementation estimates the integrals by
Monte Carlo with one million draws; that is an implementation of the
specified model, not a psychological claim that a participant samples
one million graphs or parameters.
\cite[Appendix~A, pp.~299--300]{MederMayrhoferWaldmann2014StructureInduction}
\end{hbmethod}

The likelihood distinguishes the graphs by how well their whole contingency
tables fit, rather than by the diagnostic proportion alone. If \(S_0\)
receives appreciable weight, \(g(D)\) moves toward the candidate's base
rate. The source reports about \(0.61\) for Sample~1; the other two
tables favor \(S_1\) more strongly, making the discrepancy from \(0.75\)
smaller. This number depends on the stated priors, graph space and
integrated ratio rule; it is not derivable from \(P(c\mid e)=0.75\)
alone. It is a reported model output, not a recomputation of the original
Monte Carlo program here.
\cite[fig.~7.3.1 and sec.~2.3, pp.~450--452]{MederMayrhofer2021DiagnosticCausal}

\begin{hbtheorem}{thm:p19-rat73-averaging}{proved=annex:p19.rat73-averaging}{Averaging diagnostic scores versus conditioning a mixture}
For finitely many fixed models with weights \(\rho_i\geq0\)
summing to one, put \(q_i=P_i(e)>0\) and
\(t_i=P_i(c\mid e)\). The average-of-conditionals score is
\(\sum_i\rho_it_i\). The ordinary conditional under their mixture
joint law is
\[
 \frac{\sum_i\rho_iq_it_i}{\sum_i\rho_iq_i}.
\]
These agree if all \(q_i\) are equal, but need not agree otherwise.
With equal weights, a model with \(b_c=1/2,w_c=0,w_a=1/10\)
and one with \(b_c=1/2,w_c=4/5,w_a=1/5\) give respectively
\(q_0=1/10,t_0=1/2\) and \(q_1=13/25,t_1=21/26\).
The first average is \(17/26\); the mixture conditional is \(47/62\).
\end{hbtheorem}

This is an editorial separating construction, not another source dataset.
It makes the interpretation of the source's \(g\) explicit. Its model
averages conditional scores using training-data-only weights, as in
Appendix~A6/A9. A fresh symptom treated as evidence about which model
generated that patient changes the weights as well. The latter is a
different posterior-predictive query. Calling both calculations Bayesian
does not remove the difference. Likewise, averaging only fitted parameter
means is not generally the same as averaging a nonlinear diagnostic ratio.

\section{The empirical contrast behind Figure 7.3.1}
\label{sec:p19-rat73-evidence}

The selected 2014 experiments test whether diagnostic judgments change
across data tables whose empirical diagnostic proportion stays fixed.
This can discriminate the simple data-level plug-in model from the
specified structure-induction model. It is not a test of the algebraic
validity of Bayes' rule, and the base rate fixed at \(0.5\) does not test
base-rate appreciation across prevalence levels.
\cite[pp.~278, 284--293]{MederMayrhoferWaldmann2014StructureInduction}

\subsection{Task, population and measures}

Experiment~1 recruited 36 University of G\"ottingen students,
32 women, mean age \(22.3\), for course credit or \(5\) euros.
It crossed three empirical diagnostic probabilities \(0.6,0.75,1\)
with three predictive probabilities \(0.3,0.6,0.9\), and requested
both predictive and diagnostic judgments within participant. Each of the
nine disease/substance combinations had forty individual case-symbols
on paper, randomly arranged in a five-column, eight-row grid. Participants
could inspect the table for as long as they wished; it was removed before
the two questions. Question order and data-set order were counterbalanced.
Ratings ran from \(0\) to \(7\), with endpoints expressing certainty
of absence or presence, not explicitly percentage probability.
\cite[pp.~284--285]{MederMayrhoferWaldmann2014StructureInduction}

The diagnostic ratings rose with predictive probability even within a
fixed diagnostic level. At the \(0.75\) level their means were
\(3.7,4.3,5.0\). Predictive ratings, in contrast, varied with the
predictive probability and not significantly with the diagnostic level:
the reported latter test was \(F(2,70)=1.3,p=0.28\).
This asymmetry argues against merely reversing the conditional in every
question. It does not show that every respondent explicitly represented
the two graphs. A nonsignificant effect is not proof of exact invariance.
\cite[pp.~285--287]{MederMayrhoferWaldmann2014StructureInduction}

Experiment~2 recruited 49 students from the same university, mean age
\(23.1\), \(71\%\) female, receiving course credit or \(8\) euros
per hour as part of unrelated computer experiments. Fifteen failed the
instruction comprehension test after three attempts and were excluded;
the analyzed sample was 34, mean age \(22.9\), \(76\%\) female.
Exclusion defines the population to which the reported analysis directly
applies; it is not silently a result for all recruited students.
\cite[pp.~289--290]{MederMayrhoferWaldmann2014StructureInduction}

This second experiment used the same nine frequency tables, now on a
computer in eight columns and five rows, shown for at least thirty seconds
and otherwise as long as the participant wanted. The table again disappeared
before judgment. Only diagnosis was requested, avoiding the possibility
that an accompanying predictive question suggested keeping both quantities
in mind. The target was a new person from the same population as the
learning sample: estimate disease presence given the substance in the blood.
An eleven-point scale from \(0\) to \(100\) represented percentage
probability. Data-set order was randomized.

The comprehension test expressly contrasted estimating disease presence
with three production-attribution interpretations: the disease producing
the substance, other causes producing it, or only the disease producing it.
Thus the study attempted to distinguish the instructed query rather than
assuming that the word diagnosis fixed its interpretation. The surviving
respondents' ratings still need not reveal their complete internal
representations or eliminate memory effects from removing the tables.
\cite[p.~290]{MederMayrhoferWaldmann2014StructureInduction}

\subsection{Results, rivals and uncertainty}

In Experiment~2, diagnostic probability affected judgments,
\(F(2,66)=234.4,p<0.0001\), reported \(\eta^2=0.64\).
Predictive probability also affected them across the fixed diagnostic
levels, \(F(2,66)=16.7,p<0.0001\), reported \(\eta^2=0.34\).
The interaction was not significant, \(F<1\).
These are repeated-measure comparisons of ratings, not estimates of an
effect on actual diagnosis success in a clinic. The chapter's three
Sample~1--3 judgments are a selected slice of this nine-table design.
\cite[fig.~5 and p.~290]{MederMayrhoferWaldmann2014StructureInduction}

The model comparison did not rely solely on a correlation across all
nine means. Such a correlation rewards even a model that tracks the three
diagnostic levels but misses their within-level trends. The authors also
computed the mean of three within-level correlations and the
root-mean-square deviation of model probability from the mean rating
after division by \(100\). For Experiment~2:
\[
\begin{array}{l|rrr}
 \text{candidate model}&r_{\mathrm{overall}}&r_{\mathrm{trends}}&\mathrm{RMSE}\\\hline
 \text{simple Bayes/point-fit power PC}&.859&.000&.124\\
 \text{structure induction, uniform priors}&.957&.977&.047\\
 \text{power PC, uniform priors}&.872&.282&.088\\
 \text{power PC, sparse-and-strong prior}&.856&.318&.103\\
 \text{point-fit attribution}&.920&.620&.153\\
 \text{attribution, uniform priors}&.953&.897&.135
\end{array}
\]
The \(.000\) entry is the authors' convention for a correlation undefined
because the simple model predicts no within-level variation; it is not a
measured zero correlation. The sparse-and-strong model concentrates prior
weight toward one strong productive route and a weak other route, with
the comparison using its specified prior-strength setting. It is not a
claim that causes have low base rates. The attribution models can match
some upward trends while predicting absolute scores too low for the
instructed presence question. The structure-induction model performed
best on all three listed measures in this comparison.
\cite[pp.~287--292, table~2 and Appendices~B--C]
 {MederMayrhoferWaldmann2014StructureInduction}

The authors resampled participants with replacement, taking 10,000
bootstrap samples of size 36 and 34 for the two experiments. For
Experiment~2, structure induction had the best overall correlation in
\(68.2\%\), best within-level trend correlation in \(97.3\%\),
and lowest RMSE in \(100\%\) of those resamples. These describe
relative-fit stability conditional on this empirical sample and candidate
set. They are not three population accuracies, a probability that the
mechanism is true, or results of 10,000 independent replications.
The first experiment's RMSE calculation additionally mapped the confidence
scale to \([0,1]\) by dividing by seven; that measurement interpretation
is less direct than the percentage scale in the second.
\cite[pp.~291--293, tables~2--3]{MederMayrhoferWaldmann2014StructureInduction}

The pattern supports the source's structure-sensitive account against the
particular rivals tested. It does not uniquely identify an internal
Bayesian integration process: instructions, query construal, memory,
response-scale use and other untested algorithms can matter. Two
experiments within one research program and one student recruitment
setting are not independent evidence of universality. No original
participant-level reanalysis, supplemental-code run or new replication
is reported here. Transfer to professional medicine, other cultures,
nonbinary causes, preventive relations or freely chosen diagnostic tests
requires separate evidence.

\section{Several causes, several effects and the Markov boundary}
\label{sec:p19-rat73-networks}

The source's Figure~7.3.2 compares a common effect
\(A\to B\leftarrow C\), a common cause \(A\leftarrow B\to C\),
and a chain \(A\to B\to C\). Its intended factorization is given
by the parent structure, under a Markov probability law. A causal graph
alone does not assert every pairwise dependence: the structural model
must supply suitable independence assumptions for disturbances, and
particular response functions can produce additional independences.
\cref{sec:p16-causal-factorization,sec:p16-rationality-4-2} give the
corresponding formal and modeling distinctions.
\cite[sec.~4, pp.~452--453]{MederMayrhofer2021DiagnosticCausal}

\begin{hbtheorem}{thm:p19-rat73-network}{proved=annex:p19.rat73-network}{Three factorizations and the limits of explaining away}
For finite variables and normalized probability kernels, the respective
graph factorizations are
\[
\begin{array}{l|l}
 \text{common effect}&p(a,b,c)=p(a)p(c)p(b\mid a,c)\\
 \text{common cause}&p(a,b,c)=p(b)p(a\mid b)p(c\mid b)\\
 \text{chain}&p(a,b,c)=p(a)p(b\mid a)p(c\mid b).
\end{array}
\]
The first gives marginal independence of \(A,C\); the other two
give their independence conditional on every positive-probability
value of \(B\). None forces marginal dependence in the latter two
graphs, or a strict explaining-away inequality in the first graph.
For binary independent causes with \(P(a)=r,P(c)=s\) in \((0,1)\)
and the deterministic effect \(B=A\lor C\),
\[
 P(a\mid b)=\frac{r}{r+s-rs}>r=P(a\mid b,c).
\]
By contrast, the common-effect kernel with independent fair causes and
\(P(b\mid\neg a,\neg c)=.1\),
\(P(b\mid a,\neg c)=P(b\mid\neg a,c)=.2\),
\(P(b\mid a,c)=1\) gives
\(P(a\mid b)=4/5<5/6=P(a\mid b,c)\).
It respects the same factorization and both causes increase effect
probability, but their synergy reverses the proposed discounting.
\end{hbtheorem}

The source's fever illustration uses a virus and bacterial infection as
independent, nonexclusive alternative causes. Under an appropriate
substitutable productive mechanism, discovering the bacterial cause
reduces the diagnostic evidence for the virus: the symptom has been
partly explained by another route. This is the meaning of explaining
away, not a claim that viruses and bacteria are always independent in
actual patients. The inequality needs the mechanism, not merely the
collider drawing. Deterministic OR above is a worked success case;
the synergistic kernel is a worked failure of the unrestricted claim.

For multiple effects of one cause, the source uses smoking, yellow
teeth and lung cancer as an illustrative common-cause structure.
If the likelihoods really factor conditional on both presence and absence
of that cause, diagnostic odds from two positive effects multiply their
likelihood ratios:
\[
 \frac{P(b\mid a,c)}{P(\neg b\mid a,c)}
 =\frac{P(b)}{P(\neg b)}
  \frac{P(a\mid b)}{P(a\mid\neg b)}
  \frac{P(c\mid b)}{P(c\mid\neg b)},
\]
when every displayed ratio is defined. The joint positive evidence also
must have positive probability. This is the common-cause factorization
and Bayes calculation proved in \cref{thm:p19-rat73-network}; it is
not permission to count two correlated symptoms as independent evidence.
Additional common causes, direct influence between effects, selection
or mistaken response assumptions can invalidate the product.

The chapter cites Meder and Mayrhofer's 2017 verbal-information work for
multiple-effects reasoning; it also cites Mayrhofer and Waldmann,
Park and Sloman, and Rottman and Hastie for judgments departing from
Markov restrictions. These are the authors' cited research boundaries,
not a single pooled experiment or independently qualified quantitative
record here. Likewise their mention of Stephan and Waldmann's singular
attribution work indicates a further application of attribution models,
not proof that the present student diagnosis task measures that construct.
The objections to universal causal Markov application signaled by the
chapter's Cartwright citations concern model adequacy; a participant's
departure from a stipulated kernel and a real world's departure from its
causal assumptions are different failures.
\cite[secs.~3--4, pp.~452--453]{MederMayrhofer2021DiagnosticCausal}

\section{Rational theories, empirical candidates and handbook interfaces}
\label{sec:p19-rat73-rationality}

The historical discussion contrasts early reports of conservative updating
with the heuristics-and-biases literature and later frequency-format
research. The chapter presents this as a debate about tasks, displays
and standards, not one measured population rate of rationality.
\cref{sec:p19-natural-frequencies-bayesian-reasoning} supplies the separate
frequency-format experiments. A request to count disease-positive cases
among the already observed effect-positive cases directly targets
\(\widehat P(c\mid e)\); a request to infer a new patient's disease
in an uncertain generating population can involve an additional induction
model. Improved counting does not decide which generating model is correct.
\cite[sec.~2.1, pp.~449--450]{MederMayrhofer2021DiagnosticCausal}

The conclusion of chapter~7.3 is methodological as well as normative.
Probability theory, logic-based theories, ranking theory, and bounded or
ecological rationality can offer different frameworks; even within one
probability framework the input representation, query, hypothesis space
and uncertainty treatment can differ. The authors therefore recommend
comparing explicit candidate theories rather than measuring every response
against an unexplained single yardstick. Stable departures can motivate
reconsideration of whether the imposed model is appropriate.
\cite[sec.~5, pp.~453--454]{MederMayrhofer2021DiagnosticCausal}

This argument does not establish that every stable behavior is rational,
or allow a model to be chosen after the fact solely because it fits.
A normative defense must justify its domain, priors, mechanisms, target
and losses independently; a descriptive defense must compare discriminating
predictions and actual responses. The equality of two models' point
predictions also does not establish equality of their representations.
The source's central contribution is to expose these distinctions so that
empirical disagreement is not mislabeled as an algebraic contradiction.

\begin{hbmethod}{def:p19-rat73-lookup}{A diagnostic task-selection procedure}
Specify whether the required output is a consistent fault set, a conditional
presence probability, a mechanism-specific attribution or an intervention
decision. Record the observed variables, unobserved background, sample
extent, query population, and whether numbers describe a fixed known law
or estimates from finite data. For a causal calculation, state the graph
space, generative mechanisms and independence assumptions. Distinguish
point fitting, parameter uncertainty and structure uncertainty, and state
how conditional scores are averaged. Check zero-probability evidence and
model-range failures before reporting a number.

Use an observed proportion when counting within the stated dataset is the
task; use a justified generating-model inference for a new-case task;
do not silently replace presence by production or by a counterfactual
necessity judgment. For empirical evaluation, identify the population,
learning display, comprehension exclusions, response scale, memory demands
and alternative models. A successful calculation under stipulated
assumptions and a successful explanation of human ratings are separate
achievements.
\end{hbmethod}

Consistency-based diagnosis returns the minimally withdrawn assumptions
discussed in \cref{sec:p16-consistency-based-diagnosis}; this chapter's
data tables and rated probabilities supply neither its system theory nor
its conflict family. The Bayes-net unit
\cref{sec:p16-rationality-4-2} clarifies model choice, while
\cref{sec:p16-finite-causal-models} clarifies intervention semantics.
The explanatory-preference experiments in
\cref{sec:p19-explanatory-preferences} evaluate supplied single- versus
multiple-cause explanations, not the evidence-weighted choice between
\(S_0,S_1\). Their number-of-causes simplicity must not be identified
with a marginal-likelihood preference for a two-parameter rather than
three-parameter graph. Finally \cref{sec:p19-evidence-methods} supplies
the task, population and measurement discipline used in the empirical
record above. These interfaces connect the units without treating their
different outputs as interchangeable.

\begin{proofboundary}{def:p19-rat73-boundary}{What the calculations and evidence establish}
The reciprocal annex supplies the promised probability identities,
noisy-OR derivation, mixture distinction, graph-factorization consequences,
odds calculation and separating examples. Their proofs establish formal
consequences under specified assumptions, not cognitive mechanisms.
The chapter's figure, history and methodological recommendations remain
attributed source arguments. The 2014 experiments support a bounded
descriptive comparison in the actual instructed student tasks. They do
not prove a universal theory of human rationality or a clinical method.
Other source-cited empirical applications retain their explicit cited
boundary rather than invented study denominators or effect sizes.
\end{proofboundary}
